\section{Kategorie monoidalne, symetryczne i domknięte}\label{cat:sec:monoidal}

Dotychczas iloczyn obiektów był wyróżniany przez własność uniwersalną (definicja~\ref{def:produkt}). W wielu kategoriach potrzebujemy jednak innej operacji dwuargumentowej: na przykład iloczyn tensorowy modułów nie jest na ogół ich iloczynem kategorialnym. Przejdziemy od konkretnych izomorfizmów tensorowych do aksjomatów opisujących ich wzajemną zgodność.

\subsection{Struktura monoidalna i aksjomaty koherencji}
\begin{definicja}[Kategoria monoidalna]\label{cat:def:monoidal}
Kategorią monoidalną nazywamy kategorię $\mathcal C$ z bifunktorem $\otimes:\mathcal C\times\mathcal C\to\mathcal C$, obiektem $I$ oraz izomorfizmami naturalnymi
\[
\begin{aligned}
\alpha_{A,B,C}&:(A\otimes B)\otimes C\longrightarrow A\otimes(B\otimes C),\\
\lambda_A&:I\otimes A\longrightarrow A,\\
\rho_A&:A\otimes I\longrightarrow A.
\end{aligned}
\]
Wymagamy dwóch aksjomatów koherencji. Aksjomat pięciokąta ma postać
\begin{equation}\label{cat:eq:pentagon}
\alpha_{A,B,C\otimes D}\circ\alpha_{A\otimes B,C,D}
=(\id_A\otimes\alpha_{B,C,D})\circ\alpha_{A,B\otimes C,D}
\circ(\alpha_{A,B,C}\otimes\id_D).
\end{equation}
Obie strony są morfizmami $((A\otimes B)\otimes C)\otimes D\to A\otimes(B\otimes(C\otimes D))$.
Aksjomat trójkąta ma postać
\begin{equation}\label{cat:eq:triangle-monoidal}
(\id_A\otimes\lambda_B)\circ\alpha_{A,I,B}=\rho_A\otimes\id_B,
\end{equation}
gdzie obie strony prowadzą z $(A\otimes I)\otimes B$ do $A\otimes B$.
\end{definicja}

Naturalność $\alpha$ oznacza dla $f:A\to A'$, $g:B\to B'$, $h:C\to C'$ równość
\[
\alpha_{A',B',C'}\circ((f\otimes g)\otimes h)
=(f\otimes(g\otimes h))\circ\alpha_{A,B,C}.
\]
Analogicznie $\lambda_{A'}(\id_I\otimes f)=f\lambda_A$ i $\rho_{A'}(f\otimes\id_I)=f\rho_A$. Bifunktorialność (definicja~\ref{cat:def:bifunctor}) wymaga $\id_A\otimes\id_B=\id_{A\otimes B}$ oraz
\[
(f'\otimes g')\circ(f\otimes g)=(f'\circ f)\otimes(g'\circ g)
\]
dla składanych par morfizmów.

Pięciokąt można odczytać jako dwie drogi pomiędzy pięcioma nawiasowaniami czterech czynników. Trójkąt ma następujący diagram:
\[
\begin{CD}
(A\otimes I)\otimes B @>{\alpha_{A,I,B}}>> A\otimes(I\otimes B)\\
@V{\rho_A\otimes\id_B}VV @VV{\id_A\otimes\lambda_B}V\\
A\otimes B @>{\id_{A\otimes B}}>> A\otimes B .
\end{CD}
\]
Same izomorfizmy łączności i jednostki nie wystarczają: potrzebna jest zgodność ich złożeń.

\begin{definicja}[Ścisła kategoria monoidalna]\label{cat:def:strict-monoidal}
Kategoria monoidalna jest ścisła, gdy $\otimes$ jest łączny i ma jednostkę $I$ na poziomie obiektów i morfizmów, a $\alpha,\lambda,\rho$ są identycznościami.
\end{definicja}

\begin{przyklad}[Endofunktory i składanie]\label{cat:ex:endofunctor-monoidal}
Dla małej kategorii $\mathcal C$ kategoria $[\mathcal C,\mathcal C]$ endofunktorów jest ścisła monoidalna. Przyjmujemy $F\otimes G=F\circ G$, $I=\id_{\mathcal C}$. Dla transformacji $\eta:F\Rightarrow F'$ i $\theta:G\Rightarrow G'$ iloczynem jest złożenie poziome $\eta\ast\theta:FG\Rightarrow F'G'$ z definicji~\ref{cat:def:horizontal}. Prawo wymiany (twierdzenie~\ref{cat:thm:interchange}) daje bifunktorialność, a stwierdzenie~\ref{cat:prop:horizontal-assoc} daje ścisłą łączność i jednostki. Zatem pięciokąt i trójkąt składają się z identyczności.
\end{przyklad}

\begin{twierdzenie}[Struktura kartezjańska]\label{cat:thm:cartesian-monoidal}
Każda kategoria z wybranymi iloczynami binarnymi i obiektem końcowym $1$ ma strukturę monoidalną dla $\otimes=\times$ oraz $I=1$.
\end{twierdzenie}
\begin{proof}
Bifunktor iloczynu otrzymujemy ze stwierdzenia~\ref{cat:prop:product-functor}. Zarówno $(A\times B)\times C$, jak i $A\times(B\times C)$ reprezentują iloczyn trzech obiektów: morfizm do dowolnego z nich jest jednoznacznie wyznaczony przez trzy składowe. To daje kanoniczny izomorfizm $\alpha$ zachowujący wszystkie trzy rzuty (twierdzenie~\ref{tw:produkt-unique}). Rzuty $\lambda_A:1\times A\to A$ i $\rho_A:A\times1\to A$ są izomorfizmami: odwrotności to odpowiednio $\langle !_A,\id_A\rangle$ oraz $\langle\id_A,!_A\rangle$, gdzie $!_A:A\to1$ jest jedynym morfizmem. Złożenia z rzutami pokazują, że są to odwrotności z obu stron.

Obie strony równania naturalności $\alpha$ mają po złożeniu z odpowiednimi rzutami na $A',B',C'$ składowe $f p_A$, $g p_B$, $h p_C$, gdzie $p_A,p_B,p_C$ są rzutami z obiektu wyjściowego. Jednoznaczność morfizmu do iloczynu daje naturalność. Naturalność $\lambda,\rho$ jest warunkiem zachowania odpowiedniego rzutu przez $f\times\id$ lub $\id\times f$.

W pięciokącie każde zastosowanie $\alpha$ zachowuje rzuty na czynniki, więc obie drogi mają po złożeniu z czterema końcowymi rzutami odpowiednio $p_A,p_B,p_C,p_D$. Obiekt końcowy diagramu jest iterowanym iloczynem czterech obiektów; jednoznaczność morfizmu o tych składowych daje \eqref{cat:eq:pentagon}. W trójkącie obie drogi zachowują rzuty na $A,B$, a więc są równe przez własność uniwersalną $A\times B$. Otrzymujemy \eqref{cat:eq:triangle-monoidal}.
\end{proof}

\begin{twierdzenie}[Struktura tensorowa modułów]\label{cat:thm:module-monoidal}
Niech $R$ będzie przemiennym pierścieniem z jedynką. Kategoria unitarnych $R$-modułów z iloczynem $\otimes_R$ i jednostką $R$ jest monoidalna.
\end{twierdzenie}
\begin{proof}
Wykorzystujemy konstrukcję iloczynu tensorowego, w szczególności \atlasref{AA-18::aa:tw:tensor-universal}{AA-18: własność uniwersalna}, \atlasref{AA-18::aa:tw:tensor-sym-assoc}{AA-18: łączność i symetria} oraz \atlasref{AA-18::aa:tw:tensor-unit}{AA-18: jednostka}. Izomorfizmy $\alpha,\lambda,\rho$ działają na tensorach prostych przez
\[
(m\otimes n)\otimes p\mapsto m\otimes(n\otimes p),\quad
r\otimes m\mapsto rm,\quad m\otimes r\mapsto rm .
\]
Ich dobre określenie i odwracalność wynika z przytoczonych twierdzeń. Bifunktorialność wykazano w przykładzie~\ref{cat:ex:tensor-bifunctor}.

Sprawdzimy pozostałe aksjomaty. Dla homomorfizmów $f,g,h$ obie strony naturalności $\alpha$ wysyłają $(m\otimes n)\otimes p$ na $f(m)\otimes(g(n)\otimes h(p))$. Naturalność $\lambda$ wynika z $f(rm)=rf(m)$, a dla $\rho$ działa ten sam rachunek. Tensory proste generują odpowiednie moduły, więc te równości zachodzą dla wszystkich elementów.

W pięciokącie droga z dwoma zastosowaniami $\alpha$ prowadzi przez
\[
((m\otimes n)\otimes p)\otimes q
\longmapsto (m\otimes n)\otimes(p\otimes q)
\longmapsto m\otimes(n\otimes(p\otimes q)).
\]
Droga z trzema zastosowaniami ma kolejno wartości
\[
\begin{aligned}
((m\otimes n)\otimes p)\otimes q
&\longmapsto(m\otimes(n\otimes p))\otimes q\\
&\longmapsto m\otimes((n\otimes p)\otimes q)\\
&\longmapsto m\otimes(n\otimes(p\otimes q)).
\end{aligned}
\]
Iterowane tensory proste generują moduł wyjściowy: rozwijamy każdy czynnik kolejno jako skończoną sumę tensorów prostych i korzystamy z dwuliniowości. Obie drogi są liniowe i równe na tych generatorach, więc są równe. W trójkącie pierwsza droga daje $m\otimes(rn)$, a druga $(rm)\otimes n$. Równość jest relacją zbilansowania (por. \atlasref{AA-18::aa:def:tensor-balanced}{AA-18: odwzorowania zbilansowane}). Znów wystarczy równość na generatorach.
\end{proof}

\subsection{Plecenie i symetria}
\begin{definicja}[Kategoria monoidalna z pleceniem i symetryczna]\label{cat:def:braided-symmetric}
Plecenie jest izomorfizmem naturalnym $c_{A,B}:A\otimes B\to B\otimes A$ spełniającym dwa aksjomaty sześciokąta:
\begin{align*}
\alpha_{B,C,A}\,c_{A,B\otimes C}\,\alpha_{A,B,C}
&=(\id_B\otimes c_{A,C})\,\alpha_{B,A,C}\,(c_{A,B}\otimes\id_C),\\
\alpha^{-1}_{C,A,B}\,c_{A\otimes B,C}\,\alpha^{-1}_{A,B,C}
&=(c_{A,C}\otimes\id_B)\,\alpha^{-1}_{A,C,B}\,(\id_A\otimes c_{B,C}).
\end{align*}
W pierwszym równaniu obie strony prowadzą z $(A\otimes B)\otimes C$ do $B\otimes(C\otimes A)$, w drugim z $A\otimes(B\otimes C)$ do $(C\otimes A)\otimes B$. Składamy od prawej do lewej.
Kategoria monoidalna z pleceniem jest symetryczna, jeśli dodatkowo $c_{B,A}\circ c_{A,B}=\id_{A\otimes B}$.
\end{definicja}

\begin{stwierdzenie}[Symetria w przykładach kartezjańskich i tensorowych]\label{cat:prop:cartesian-tensor-symmetry}
Struktury z twierdzeń~\ref{cat:thm:cartesian-monoidal} i \ref{cat:thm:module-monoidal} są symetryczne.
\end{stwierdzenie}
\begin{proof}
W przypadku kartezjańskim bierzemy $c_{A,B}=\langle p_B,p_A\rangle$. Dwa złożenia zamian zachowują oba rzuty, więc są identycznościami. Naturalność wynika z tego, że obie drogi dla $f:A\to A'$ i $g:B\to B'$ mają składowe $g p_B$ oraz $f p_A$. Pierwszy sześciokąt daje po obu drogach uporządkowane składowe $B,C,A$, drugi $C,A,B$, z zachowaniem rzutów na poszczególne czynniki. Jednoznaczność morfizmów do iloczynu dowodzi obu równości.

W przypadku modułów stosujemy izomorfizm $c_{M,N}(m\otimes n)=n\otimes m$ z \atlasref{AA-18::aa:tw:tensor-sym-assoc}{AA-18}. Naturalność sprawdzamy na tensorach prostych: obie drogi dają $g(n)\otimes f(m)$. W pierwszym sześciokącie obie drogi wysyłają $(m\otimes n)\otimes p$ na $n\otimes(p\otimes m)$; w drugim obie wysyłają $m\otimes(n\otimes p)$ na $(p\otimes m)\otimes n$. Podwójna zamiana daje $m\otimes n$. Liniowość i generowanie przez tensory proste kończą dowód.
\end{proof}

\begin{przyklad}[Symetria przestrzeni z gradacją]\label{cat:ex:graded-symmetry}
Dla $\mathbb Z$-stopniowanych przestrzeni wektorowych nad ciałem $k$ i odwzorowań zachowujących stopień rozważamy standardowy iloczyn tensorowy z gradacją $(V\otimes W)_n=\bigoplus_{p+q=n}V_p\otimes W_q$. Jednostką jest $k$ w stopniu zero. Łączność i jednostki są odziedziczone z iloczynu tensorowego, zachowują stopień, a ich aksjomaty wynikają z twierdzenia~\ref{cat:thm:module-monoidal}.
Na jednorodnych wektorach $v\in V_p,w\in W_q$ definiujemy
\[
c(v\otimes w)=(-1)^{pq}w\otimes v.
\]
Wzór na każdej parze składowych stopniowanych jest dwuliniowy i daje liniowe odwzorowanie z sumy prostej; zachowuje stopień $p+q$. Naturalność wynika z zachowywania stopni przez morfizmy. Dwukrotna zamiana daje czynnik $(-1)^{2pq}=1$, więc $c$ jest izomorfizmem i spełnia symetrię. Dla trzeciego jednorodnego wektora stopnia $r$ pierwszy sześciokąt daje znaki $(-1)^{p(q+r)}$ i $(-1)^{pq+pr}$, drugi $(-1)^{(p+q)r}$ i $(-1)^{pr+qr}$. Znaki są równe, a kolejność tensorów po obu drogach taka sama. Jednorodne tensory proste generują rozpatrywane moduły, więc oba sześciokąty zachodzą. W charakterystyce $2$ wszystkie te znaki są równe $1$.
\end{przyklad}

\subsection{Domkniętość monoidalna}
\begin{definicja}[Prawa domkniętość monoidalna]\label{cat:def:monoidal-closed}
Kategoria monoidalna jest prawostronnie domknięta, gdy dla każdego $A$ funktor $-\otimes A$ ma prawe sprzężenie $[A,-]$. Oznacza to naturalną bijekcję
\[
\Hom_{\mathcal C}(X\otimes A,B)\cong\Hom_{\mathcal C}(X,[A,B]).
\]
Obiekt $[A,B]$ nazywamy wewnętrznym obiektem Hom. W kategorii symetrycznej domkniętość z jednej strony pociąga domkniętość z drugiej: symetria daje naturalny izomorfizm $A\otimes-\cong-\otimes A$, więc powyższe bijekcje można złożyć z bijekcjami indukowanymi przez ten izomorfizm.
\end{definicja}

\begin{przyklad}[Domkniętość kartezjańska i tensorowa]\label{cat:ex:closed-monoidal}
W kategorii kartezjańsko domkniętej $[A,B]=B^A$ (definicja~\ref{def:exponential}, twierdzenie~\ref{cat:thm:exp-naturality}). W kategorii modułów nad przemiennym pierścieniem $R$ przyjmujemy $[A,B]=\operatorname{Hom}_R(A,B)$ z punktowymi działaniami modułu. Naturalną bijekcję wyznaczają wzory
\[
\Phi(f)(x)(a)=f(x\otimes a),\qquad
\Psi(g)(x\otimes a)=g(x)(a).
\]
Ich dobre określenie i wzajemną odwrotność daje \atlasref{AA-18::aa:tw:tensor-hom-adjunction}{AA-18: sprzężenie tensor--Hom}. Dla $u:X'\to X$, $v:B\to B'$ naturalność sprawdzamy bezpośrednio:
\[
\Phi(vf(u\otimes\id_A))(x')(a)=v(f(u(x')\otimes a)).
\]
Jest to dokładnie wartość złożenia $\operatorname{Hom}_R(A,v)\circ\Phi(f)\circ u$ w $x'$ i $a$. W ten sposób uzyskujemy sprzężenie, nie tylko bijekcję poszczególnych zbiorów. Nie potrzeba skończonej wymiarowości ani skończonego generowania modułów.
\end{przyklad}

\begin{stwierdzenie}[Zachowywanie kogranic przez iloczyn monoidalny]\label{cat:prop:closed-colimits}
W prawostronnie domkniętej kategorii monoidalnej funktor $-\otimes A$ zachowuje istniejące małe kogranice. W przypadku symetrycznym dotyczy to także $A\otimes-$.
\end{stwierdzenie}
\begin{proof}
Z definicji~\ref{cat:def:monoidal-closed} jest to lewe sprzężenie. Stosujemy twierdzenie~\ref{tw:adjoints-preserve}. W przypadku symetrycznym drugi funktor jest naturalnie izomorficzny z pierwszym; izomorfizmy przenoszą własność uniwersalną kogranicy.
\end{proof}

\begin{uwaga}
Koherencja dla dowolnie długich wyrażeń monoidalnych, funktory monoidalne, obiekty monoidalne oraz obiekty posiadające dualność wymagają dalszych definicji i twierdzeń. Powyższe przykłady dowodzą konkretnych aksjomatów pięciokąta i trójkąta, a nie ogólnego twierdzenia Mac Lane'a o koherencji. W dalszej części omówimy je oddzielnie.
\end{uwaga}
